% Indicateurs de fiabilité. Une seule série, donc une seule teinte et pas de
% légende : le titre de l'axe suffit à la nommer. Le rail clair rappelle que
% l'échelle va jusqu'à 100 %, ce qu'une barre nue laisse deviner.
% Indicateurs de fiabilité, en pourcentage
% Généré par rapport/build_data.py — NE PAS ÉDITER.
\newcommand{\FiabiliteLabels}{{Citation systématique},{Citations vérifiées},{Article attendu cité},{Abstention hors périmètre}}
\newcommand{\FiabiliteCoords}{(96.1,0) (81.1,1) (88.2,2) (100.0,3)}
\newcommand{\FiabiliteDetailA}{49/51}
\newcommand{\FiabiliteDetailB}{77/95}
\newcommand{\FiabiliteDetailC}{45/51}
\newcommand{\FiabiliteDetailD}{9/9}

\begin{tikzpicture}[
  detail/.style={font=\sffamily\scriptsize, text=slate, anchor=west}]
\begin{axis}[
  barresH,
  height=4.8cm,
  % Largeur réduite : les libellés à gauche et les effectifs à droite
  % s'ajoutent à l'aire tracée.
  width=0.62\linewidth,
  xmin=0, xmax=134,
  xtick={0,25,50,75,100},
  xticklabel={\pgfmathprintnumber{\tick}\,\%},
  ytick={0,...,3},
  % Libellés en clair : pgfplots ne développe pas une macro donnée à
  % « yticklabels ». L'ordre suit celui de la liste « fiab » dans
  % rapport/build_data.py, qui fixe l'ordre des barres.
  yticklabels={{Citation systématique},{Citations vérifiées},
               {Article attendu cité},{Abstention hors périmètre}},
  yticklabel style={font=\sffamily\scriptsize, text=ink, align=right},
  xlabel={part des cas conformes},
  bar width=11pt,
  every axis plot/.append style={bar shift=0},
]
% Rail de fond : la portion non atteinte.
\addplot[draw=none, fill=mist, forget plot]
  coordinates {(100,0) (100,1) (100,2) (100,3)};
\addplot[draw=inptblue, fill=inptblue, fill opacity=0.9,
         nodes near coords={\pgfmathprintnumber[precision=0]{\pgfplotspointmeta}\,\%},
         nodes near coords style={font=\sffamily\scriptsize\bfseries, text=ink,
                                  anchor=west, xshift=3pt}]
  coordinates {\FiabiliteCoords};

% Effectif réel derrière chaque pourcentage : 93 % sur 45 citations n'a pas
% le même poids que 100 % sur 5 questions, et la figure doit le dire.
\node[detail] at (axis cs:116,0) {\FiabiliteDetailA};
\node[detail] at (axis cs:116,1) {\FiabiliteDetailB};
\node[detail] at (axis cs:116,2) {\FiabiliteDetailC};
\node[detail] at (axis cs:116,3) {\FiabiliteDetailD};
\end{axis}
\end{tikzpicture}
